% !TeX program = xelatex
% Compile twice with XeLaTeX (or pdfLaTeX).
\documentclass[a4paper,11pt]{article}
\usepackage[margin=25mm]{geometry}
\usepackage[T1]{fontenc}
\usepackage{lmodern}
\usepackage{amsmath,amssymb,amsthm,mathtools}
\usepackage{microtype}
\usepackage[colorlinks=true,linkcolor=blue,citecolor=blue,urlcolor=blue]{hyperref}
\newtheorem{proposition}{Proposition}
\newtheorem{lemma}{Lemma}
\theoremstyle{remark}
\newtheorem*{remark}{Remark}
\DeclareMathOperator{\sfpart}{sf}
\DeclareMathOperator{\TV}{TV}
\newcommand{\Rset}{\mathcal R}
\newcommand{\PrD}{\mathbb P_{\Delta}}
\newcommand{\PrU}{\mathbb P_{\mathrm U}}
\newcommand{\N}{\mathbb N}
\hypersetup{
  pdftitle={A congruence obstruction to Heuristic 1 in POKE-2D},
  pdfauthor={Wenbiao Chen}
}

\begin{document}

\noindent Title: \textbf{A congruence obstruction to Heuristic 1 in POK\'E-2D}
\par\vspace{1em}
\noindent
\begin{minipage}[t]{0.76\textwidth}
Author: Wenbiao Chen, Institute of Information Engineering,\\
Chinese Academy of Sciences, China.
\end{minipage}\hfill
\begin{minipage}[t]{0.23\textwidth}
\raggedleft Date: Oct. 14, 2026
\end{minipage}
\par\vspace{1.4em}

This note concerns the probability estimates in \emph{On weak keys of
POK\'E}~\cite{Moriya2026}. For POK\'E-2D, the square-free part of
$D'^2-\delta$ satisfies congruence restrictions absent from the uniform
reference model in Heuristic~1. We give an explicit statistical-distance
bound and an admissible candidate set with zero actual weak-key probability
but a constant-order estimate in the paper's regime $D'^2\asymp\delta_m$.
For the actual KeyGen parameters, the square condition targeted by
Algorithm~1 is impossible. Detailed proofs are in the appendices.

\section{Setting and the distributional obstruction}

We use the paper's notation:
\[
 \Delta=\{q(A-q):q\in\{1,\ldots,A-1\}\},
 \qquad \delta_m=\max\Delta=\frac{A^2}{4}.
\]
Assume $8\mid A$, which holds in the original POK\'E-2D
parameterization~\cite{BassoMaino2025} with $A=2^a$ and $a\ge3$.
We retain the paper's conditions
$\gcd(A,D)=1$, $D'\mid D$, and $D'\ge\sqrt{\delta_m}$.
Thus $D'$ is odd and $D'^2>\delta_m$.

Let $\sfpart(n)$ denote the unique square-free integer $r$ such that
$n=rc^2$ for a positive integer $c$. Let
\[
 \delta\overset{\$}{\leftarrow}\Delta,
 \qquad
 \delta_0\overset{\$}{\leftarrow}\{1,\ldots,\delta_m\},
 \qquad
 S=\sfpart(D'^2-\delta),\quad S_0=\sfpart(D'^2-\delta_0).
\]
Write $\PrD$ and $\PrU$ for these two probability models.
The first sampling rule is uniform over the distinct degrees in $\Delta$,
as in the paper. For square-free $r$, Heuristic~1 asserts the unquantified approximation
\begin{equation}\label{eq:H1}
 \PrD[S=r]\approx\PrU[S_0=r].
\end{equation}

Here $\mathcal L(S)$ denotes the probability distribution of $S$.
The total variation distance is
\begin{align*}
 \TV\bigl(\mathcal L(S),\mathcal L(S_0)\bigr)
 &:=\sup_{\mathcal E\subseteq\N}
       |\PrD[S\in\mathcal E]-\PrU[S_0\in\mathcal E]|\\
 &=\frac12\sum_{r\ge1}|\PrD[S=r]-\PrU[S_0=r]|.
\end{align*}

\begin{proposition}[Explicit distributional discrepancy]\label{prop:gap}
Let $\varepsilon=(\lfloor\log_4\delta_m\rfloor+2)/\delta_m$.
The degree set $\Delta$ imposes the restriction
\[
 S\bmod8\in\{1,2,5\}.
\]
In contrast, for each $u\in\{1,2,3,5,6,7\}$,
\[
 \left|\PrU[S_0\equiv u\pmod8]-\frac16\right|
 \le\varepsilon.
\]
Consequently, for the event
$\mathcal B=\{r\in\N:r\bmod8\in\{3,6,7\}\}$,
\begin{equation}\label{eq:forbidden}
 \PrD[S\in\mathcal B]=0,
 \qquad
 \left|\PrU[S_0\in\mathcal B]-\frac12\right|\le3\varepsilon,
\end{equation}
and hence
\begin{equation}\label{eq:TV}
 \TV\bigl(\mathcal L(S),\mathcal L(S_0)\bigr)
 \ge\frac12-3\varepsilon.
\end{equation}
\end{proposition}

Since $\varepsilon\to0$ as $\delta_m\to\infty$, Heuristic~1 cannot hold
in total variation with vanishing error. If interpreted only as an additive $o(1)$
error for each fixed $r$, it is not contradicted by this argument.
However, that interpretation neither controls the relative error in the
small probabilities estimated in Theorem~14 nor justifies summing the
errors over a growing set $\Rset$ in Theorem~17.

\section{An admissible candidate set for Theorem 17}

Definition~16 requires only that $\Rset\subseteq\{1,\ldots,\delta_m\}$
consist of square-free integers. In particular, it permits
\begin{equation}\label{eq:R}
 \Rset=\{1\le r\le\delta_m:r\text{ is square-free},\
                                  r\equiv3\pmod4\}.
\end{equation}
This is a parameter-dependent family of sets. The restriction to the
smallest square-free integers is introduced separately in Section~4.2 of
the paper and is not a hypothesis of Theorem~17.

Denote the probability expression displayed in Theorem~17 by
\begin{equation}\label{eq:F}
 F(\Rset)=
 \frac{1}{D'+\sqrt{D'^2-\delta_m}}
 \sum_{r\in\Rset}\frac1{\sqrt r}.
\end{equation}

\begin{proposition}[An allowed set with zero actual probability]\label{prop:R}
For the set~\eqref{eq:R}, under any degree distribution supported on
$\Delta$,
\begin{equation}\label{eq:zero}
 \PrD[\exists r\in\Rset,\ c\in\N:\ D'^2-\delta=rc^2]=0.
\end{equation}
Nevertheless,
\begin{equation}\label{eq:Fasym}
F(\Rset)
=
\frac{4}{\pi^2}\,
\frac{\sqrt{\delta_m}}{D'+\sqrt{D'^2-\delta_m}}
+
O\!\left(
\frac{\log\delta_m}{D'+\sqrt{D'^2-\delta_m}}
\right).
\end{equation}
Thus, in the regime $D'^2\asymp\delta_m$, $F(\Rset)=\Theta(1)$.
In particular, if $D'^2/\delta_m\to1$, then
$F(\Rset)\to4/\pi^2\approx0.4053$.
\end{proposition}

In that regime, the exact uniform reference probability for the same set
is also bounded away from zero; see Appendix~\ref{app:counter}.
Thus the failure of the reference model is not merely an artefact of
approximating integer counts.

The expression in Theorem~17 therefore cannot approximate the actual
POK\'E-2D probability with additive error $o(1)$ uniformly over all its
allowed candidate sets. This identifies a limitation in applying the
heuristic estimate; it does not refute a conditional implication whose
premise is Heuristic~1. Moreover, $\#\Rset=\Theta(\delta_m)$, so this example
does not invalidate the conditional recovery algorithm or establish an
attack contradicting the proposed security parameters. The actual
probability being smaller also leaves a one-sided upper bound unrefuted.

\section{Effect of the KeyGen coprimality conditions}

\begin{remark}
In the original POK\'E-2D parameterization~\cite{BassoMaino2025},
valid secret degrees additionally satisfy $\gcd(\delta,AD)=1$.
Since $A$ is even, this forces $q$ to be odd. Hence
\[
 \delta=q(A-q)\equiv7\pmod8,\qquad
 D'^2-\delta\equiv2\pmod8,\qquad
 \sfpart(D'^2-\delta)\equiv2\pmod8.
\]
For any sampling weights on these admissible degrees, the event
$\sfpart(D'^2-\delta)\not\equiv2\pmod8$ has probability zero.
Keeping the paper's unconditioned reference distribution unchanged,
its probability differs from $5/6$ by at most $\varepsilon$.
Thus the statistical distance lower bound strengthens to
\[
 \frac56-\varepsilon.
\]
The zero-probability statement for~\eqref{eq:R} remains valid.
\end{remark}

In particular, Algorithm~1 of~\cite{Moriya2026}, based on Theorem~7,
targets the condition
\[
 D'^2-\deg\varphi=c^2.
\]
This would require $\sfpart(D'^2-\deg\varphi)=1$, contradicting the
congruence above. Hence no valid POK\'E-2D secret key satisfies this
condition for any admissible $D'\mid D$: the square-condition weak-key
class targeted by Algorithm~1 has probability zero.
This does not exclude the generalized case in Algorithm~2 with
$r\equiv2\pmod8$.

A revised probability model should account for these local restrictions.
The approximation in Heuristic~1 also needs an explicit error criterion
and a range of candidate sets over which errors may be summed.
Accounting for these congruence restrictions is necessary, but does not
by itself justify the probability estimates.

\appendix
\section{Proof of the distributional discrepancy}\label{app:distribution}

For $\delta=q(A-q)$, odd $D'$ and $8\mid A$ give
\[
 D'^2-\delta\equiv1+q^2\pmod8\in\{1,2,5\}.
\]
In particular, $4\nmid D'^2-\delta$. Writing $D'^2-\delta=rc^2$
with $r$ square-free, $c$ must be odd; therefore $c^2\equiv1\pmod8$
and $r\equiv D'^2-\delta\pmod8$.

For the reference model, every positive integer $n$ has a unique
decomposition $n=4^j m$ with $4\nmid m$. Since the square factor remaining
in $m$ is odd, $\sfpart(n)\equiv m\pmod8$. Consequently, for
$u\in\{1,2,3,5,6,7\}$,
\begin{equation}\label{eq:progressions}
 \sfpart(n)\equiv u\pmod8
 \quad\Longleftrightarrow\quad
 n=4^j(8k+u),\qquad j,k\ge0.
\end{equation}
For fixed $u$, the sets indexed by $j$ are disjoint.

Now $D'^2-\delta_0$ is uniform on the $\delta_m$ integers in
\[
 I=\{D'^2-\delta_m,\ldots,D'^2-1\}.
\]
Set $M=D'^2-1$ and $J=\lfloor\log_4\delta_m\rfloor$.
For $0\le j\le J$, let $h=8\cdot4^j$ and $a=4^j u$.
The exact count is
\begin{align}
 \#\{n\in I:n\equiv a\pmod h\}
 &=\left\lfloor\frac{M-a}{h}\right\rfloor
   -\left\lfloor\frac{M-\delta_m-a}{h}\right\rfloor\notag\\
 &=\frac{\delta_m}{h}+e_{j,u},\qquad |e_{j,u}|<1.
 \label{eq:floor}
\end{align}
Indeed, $e_{j,u}$ is the difference of the fractional parts of
$(M-\delta_m-a)/h$ and $(M-a)/h$. No divisibility assumption is needed.
All remaining terms, with $j>J$, are divisible by $4^{J+1}$.
Since $4^{J+1}>\delta_m$, at most one integer in $I$ can contribute
to this tail. Denote its contribution for residue $u$ by $T_u\in\{0,1\}$.
It follows that
\[
 \PrU[S_0\equiv u\pmod8]
 =\frac16(1-4^{-(J+1)})
   +\frac{\sum_{j=0}^{J}e_{j,u}+T_u}{\delta_m}.
\]
Now $0<\delta_m/(6\cdot4^{J+1})<1/6$, so
\[
 \left|T_u-\frac{\delta_m}{6\cdot4^{J+1}}\right|<1.
\]
Combining this with $|e_{j,u}|<1$ gives
\[
 \left|\PrU[S_0\equiv u\pmod8]-\frac16\right|
 \le\frac{J+2}{\delta_m}=\varepsilon.
\]
This bound is independent of the position of $I$ and requires no upper
bound on $D'^2/\delta_m$. Summing the classes $3,6,7$ gives
\eqref{eq:forbidden}. Finally, total variation is the supremum of the
probability difference over events, giving~\eqref{eq:TV}.

\section{Square-free counting and the weighted sum}\label{app:counting}

\begin{lemma}\label{lem:count}
For $x\ge2$, let
\[
 N(x)=\#\{n\le x:n\text{ is square-free},\ n\equiv3\pmod4\}.
\]
Then
\begin{equation}\label{eq:count}
 N(x)=\frac2{\pi^2}x+O(\sqrt x),
 \qquad
 \sum_{\substack{n\le x\\n\text{ square-free}\\n\equiv3\ (4)}}
 \frac1{\sqrt n}
 =\frac4{\pi^2}\sqrt x+O(\log x).
\end{equation}
\end{lemma}

\begin{proof}
Let $\mu$ be the M\"obius function. The identity
\[
 \mu(n)^2=\sum_{a^2\mid n}\mu(a)
\]
follows by expanding $\prod_{p^2\mid n}(1-1)$: the empty product is $1$,
whereas the product is $0$ if a squared prime divides $n$.
Interchanging finite sums gives
\[
 N(x)=
 \sum_{a\le\sqrt x}\mu(a)
 \#\{b\le x/a^2:a^2b\equiv3\pmod4\}.
\]
Even $a$ contribute nothing. For odd $a$, $a^2\equiv1\pmod4$, and
\[
 \#\{b\le y:b\equiv3\pmod4\}
 =\left\lfloor\frac{y+1}{4}\right\rfloor
 =\frac y4+O(1).
\]
Consequently,
\begin{equation}\label{eq:mobiuscount}
 N(x)=\frac x4\sum_{\substack{a\le\sqrt x\\a\ {\rm odd}}}
                      \frac{\mu(a)}{a^2}+O(\sqrt x).
\end{equation}
The error follows from $|\mu(a)|\le1$ and at most $\sqrt x$ summands.
Also,
\[
 \sum_{a>\sqrt x}\frac{|\mu(a)|}{a^2}=O(x^{-1/2}),
\]
so completing the sum in~\eqref{eq:mobiuscount} to infinity changes the
answer by $O(\sqrt x)$.

For completeness, the relevant Euler product follows from
\[
 \prod_{\substack{3\le p\le y\\p\ {\rm prime}}}(1-p^{-2})
 =
 \sum_{\substack{a\ge1,\ a\ {\rm odd}\\p\mid a\Rightarrow p\le y}}
 \frac{\mu(a)}{a^2}.
\]
Expanding the finite product selects distinct odd primes, and their
product $a$ has coefficient $\mu(a)$. The omitted terms have $a>y$,
with absolute sum at most $\sum_{a>y}a^{-2}\to0$. Hence
\[
 \sum_{\substack{a\ge1\\a\ {\rm odd}}}\frac{\mu(a)}{a^2}
 =\prod_{p\ne2}(1-p^{-2})
 =\frac{\zeta(2)^{-1}}{1-2^{-2}}
 =\frac8{\pi^2}.
\]
Substituting into~\eqref{eq:mobiuscount} proves the first formula.

To prove the weighted formula, let
$a_n=\mathbf1_{\{n\ {\rm square\text{-}free},\,n\equiv3\ (4)\}}$.
Then $N(t)=\sum_{n\le t}a_n$. For $f\in C^1([1,x])$, the identity
$f(n)=f(x)-\int_n^x f'(t)\,dt$ and linearity of integration over
a finite sum imply
\begin{align*}
 \sum_{n\le x}a_nf(n)
 &=N(x)f(x)
   -\int_1^x\left(\sum_{n\le x}a_n\mathbf1_{\{n\le t\}}\right)
     f'(t)\,dt\\
 &=N(x)f(x)-\int_1^x N(t)f'(t)\,dt.
\end{align*}
Applying this exact identity to $f(t)=t^{-1/2}$ yields
\[
 \sum_{n\le x}\frac{a_n}{\sqrt n}
 =\frac{N(x)}{\sqrt x}
  +\frac12\int_1^x N(t)t^{-3/2}\,dt.
\]
Write $N(t)=\alpha t+E(t)$, where
$\alpha=2/\pi^2$ and $E(t)=O(\sqrt t)$ for $t\ge1$.
The endpoint contributes $\alpha\sqrt x+O(1)$, and the main term in
the integral contributes
\[
 \frac{\alpha}{2}\int_1^x t^{-1/2}\,dt
 =\alpha(\sqrt x-1).
\]
The remaining integral is bounded in absolute value by
\[
 \frac12\int_1^x |E(t)|t^{-3/2}\,dt
 =O\!\left(\int_1^x\frac{dt}{t}\right)=O(\log x).
\]
Combining these terms proves the second formula.
\end{proof}

\section{The candidate-set counterexample and reference counts}
\label{app:counter}

For $\delta\in\Delta$,
\[
 D'^2-\delta\equiv1+q^2\pmod4\in\{1,2\}.
\]
For $r\equiv3\pmod4$, however, $rc^2$ is $0$ or $3$ modulo $4$.
Thus the equality $D'^2-\delta=rc^2$ is impossible, proving
\eqref{eq:zero} for any sampling weights on $\Delta$.

By Lemma~\ref{lem:count},
\[
 \#\Rset=\frac2{\pi^2}\delta_m+O(\sqrt{\delta_m}),
 \qquad
 \sum_{r\in\Rset}\frac1{\sqrt r}
 =\frac4{\pi^2}\sqrt{\delta_m}+O(\log\delta_m).
\]
Substitution in~\eqref{eq:F} and division of numerator and denominator
by $\sqrt{\delta_m}$ prove~\eqref{eq:Fasym}.
If $D'^2\asymp\delta_m$, the denominator in~\eqref{eq:Fasym} is
bounded above and bounded below by $1$, whereas the numerator tends to
$4/\pi^2$. This proves the constant-order assertion.
If $D'^2/\delta_m\to1$, the denominator tends to $1$.

\paragraph{Exact reference probability.}
Let $b=\lceil D'/\sqrt{\delta_m}\rceil$ and count only reference degrees with
$D'^2-\delta_0=b^2r$, $r\in\Rset$. Such $r$ lie in
\[
 \left[\frac{D'^2-\delta_m}{b^2},
       \frac{D'^2-1}{b^2}\right]\subset(0,\delta_m].
\]
Every square-free $r\equiv3\pmod4$ in this interval yields a distinct
valid $\delta_0$ and satisfies $\sfpart(b^2r)=r$.
Lemma~\ref{lem:count}, with an $O(1)$ endpoint correction, gives
\[
 \#\{r\text{ in this interval}:r\text{ square-free},\
                              r\equiv3\pmod4\}
 =\frac{2(\delta_m-1)}{\pi^2b^2}+O(\sqrt{\delta_m}).
\]
For endpoints below $2$, the same estimate follows by absorbing a bounded
error. Dividing by $\delta_m$ proves
\begin{equation}\label{eq:exactreference}
 \PrU[S_0\in\Rset]\ge
 \frac2{\pi^2b^2}-O(\delta_m^{-1/2}).
\end{equation}
Thus the exact reference distribution already assigns positive constant
mass, when $D'^2\asymp\delta_m$ and hence $b=O(1)$, to a set that is
impossible under the actual degree support.

\paragraph{Rounding and accumulation of errors.}
For each square-free $r$, the exact reference probability is
\begin{align*}
 \PrU[S_0=r]
 &=\frac1{\delta_m}
 \left(
  \left\lceil\frac{D'}{\sqrt r}\right\rceil
  -\left\lceil\sqrt{\frac{D'^2-\delta_m}{r}}\right\rceil
 \right)\\
 &=\frac1{\sqrt r\bigl(D'+\sqrt{D'^2-\delta_m}\bigr)}
   +\eta_r,\qquad |\eta_r|\le\frac1{\delta_m}.
\end{align*}
Define the model error
$\varepsilon_r=\PrD[S=r]-\PrU[S_0=r]$. For any admissible $\Rset$,
\begin{equation}\label{eq:errors}
 \PrD[S\in\Rset]
 =F(\Rset)+\sum_{r\in\Rset}\varepsilon_r
           +\sum_{r\in\Rset}\eta_r,
 \qquad
 \left|\sum_{r\in\Rset}\eta_r\right|
 \le\frac{\#\Rset}{\delta_m}.
\end{equation}
Pointwise additive $o(1)$ bounds on $\varepsilon_r$ do not bound their
sum over growing sets. For the family~\eqref{eq:R}, the available
rounding bound is also constant-sized. Statements~\eqref{eq:forbidden}
and~\eqref{eq:exactreference} establish the model discrepancy independently
of these rounding errors.

\begin{thebibliography}{2}
\bibitem{Moriya2026}
Tomoki Moriya.
\newblock \emph{On weak keys of POK\'E}.
\newblock Cryptology ePrint Archive, Paper 2026/1002,
version dated 2026-05-21.
\newblock \url{https://eprint.iacr.org/2026/1002}.
\bibitem{BassoMaino2025}
Andrea Basso and Luciano Maino.
\newblock \emph{POK\'E: A Compact and Efficient PKE from
Higher-dimensional Isogenies}.
\newblock EUROCRYPT 2025, pp.~94--123.
\newblock \url{https://eprint.iacr.org/2024/624}.
\end{thebibliography}
\end{document}

